\section{The Heat Equation} \label{sec.heat_equation}
The heat equation is a fundamental partial differential equation that describes how heat diffuses through a given region over time. It is given by the following equation:
\begin{equation}\label{eqn.heat_eqn}
    u_t = a^2 u_{xx}
\end{equation}
This is the one-dimensional heat equation, generally modeling a rod,  where $u(x,t)$ represents the temperature at position $x$ and time $t$, and $a^2$ is the thermal diffusivity constant. The subscript notation denotes partial derivatives, with $u_t$ representing the partial derivative of $u$ with respect to time, and $u_{xx}$ representing the second partial derivative of $u$ with respect to space.
It is usually presented with initial and boundary conditions, which for the following analysis we will use the following:
\begin{align*} \label{ic.heat_eqn}
    u(x,0) &= f(x) \quad \text{for } 0 < x < L \\
    u(0,t) &= 0 \quad \text{for } t > 0 \\
    u(L,t) &= 0 \quad \text{for } t > 0
\end{align*}

The initial condition $u(x,0) = f(x)$ specifies the temperature distribution along the rod at time $t=0$, while the boundary conditions $u(0,t) = 0$ and $u(L,t) = 0$ indicate that the ends of the rod are held at zero temperature for all times $t > 0$. The goal is to find a function $u(x,t)$ that satisfies the heat equation \eqref{eqn.heat_eqn} along with the given initial and boundary conditions. This problem can be solved using various methods, including separation of variables, Fourier series, and numerical techniques such as finite difference methods.

\subsection{Separation of Variables}\label{subsec.separation_of_variables}
This is a common technique for solving partial differential equations and the process will be shown for the heat equation \eqref{eqn.heat_eqn}. The idea is to assume that the solution can be written as a product of two functions, one depending only on $x$ and the other depending only on $t$. Suppose:
\[
    u(x,t) = X(x)T(t)
\]
where $X(x)$ is a function of $x$ alone and $T(t)$ is a function of $t$ alone. We make no claim about what $X$ and $T$ are yet. Let's transform the heat equation into separation form:
\[
\begin{cases}
    u_t = X(x)T'(t) \\
    u_x = X'(x)T(t) \\
    u_{xx} = X''(x)T(t)
\end{cases}
\implies X(x)T'(t) = a^2 X''(x)T(t)
\]
This doesn't look like it's helpful yet, but we can rearrange it to get:
\[
\frac{T'(t)}{a^2 T(t)} = \frac{X''(x)}{X(x)}
\]
so now each side only has one variable. Since our equation must hold for all $x$ and $t$, we can set both sides equal to a constant, say $k$: 
\[
\frac{T'(t)}{a^2 T(t)} = \frac{X''(x)}{X(x)} = k
\]
This gives us two ordinary differential equations:
\[
\begin{cases}
    T'(t) - a^2 k T(t) = 0 \\
    X''(x) - k X(x) = 0
\end{cases}
\]
These are much easier to solve. 
\subsubsection{What can $k$ be?}\label{subsubsec.k_values}
The value of the constant $k$ at this point is not known. Let's try to refine what $k$ can be. $k$ can be either of the following: $k >0$, $k=0$, or $k<0$. \\
\textit{Case 1: $k = 0$:}
In this case, the equation for $X(x)$ becomes:
\[
\begin{cases}
    T'(t) = 0 \\
    X''(x) = 0
\end{cases}
\implies 
    X(x) = C_1 x + C_2
\]
for some constants $C_1$ and $C_2$. Let's check the initial and boundary conditions:
\[
u(0,t) = X(0)T(t) = C_2 \cdot T(t) = 0 \\
\]
For this to hold for all $t$, we must have $C_2 = 0$. Now we check the other boundary condition:
\[
u(L,t) = X(L)T(t) = C_1 L \cdot T(t) = 0
\]
For this to hold for all $t$, we must have $C_1 = 0$. Thus, the only solution for $k=0$ is the trivial solution $u(x,t) = 0$ which we are not interested in.\\
\textit{Case 2: $k > 0$:}
From ODEs:
\[
X''(x) - k X(x) = 0
\implies 
X(x) = C_1 e^{\sqrt{k} x} + C_2 e^{-\sqrt{k} x}
\]
Check initial and boundary conditions, $u(0,t)=0$:
\[
u(0,t) = X(0)T(t) = (C_1 + C_2)T(t) = 0
\implies C_2 = -C_1 \implies
X(x) = C_1 (e^{\sqrt{k} x} - e^{-\sqrt{k} x}) 
\]
Check second boundary condition, $u(L,t)=0$:
\[
u(L,t) = X(L)T(t) = C_1 (e^{\sqrt{k} L} - e^{-\sqrt{k} L}) T(t) = 0
\]
Since $e^{\sqrt{k} L} \geq e^{-\sqrt{k} L}$ for $k > 0$, $e^{\sqrt{k} L} - e^{-\sqrt{k} L} \neq 0$, so we must have $C_1 = 0$. Thus, the only solution for $k > 0$ is also the trivial solution $u(x,t) = 0$.\\
\textit{Case 3: $k < 0$:}
Let $k = -\lambda^2$ for some $\lambda > 0$. Then the equation for $X(x)$ becomes:
\[
X''(x) + \lambda^2 X(x) = 0
\]
This is a second-order linear homogeneous ODE with constant coefficients. The characteristic equation is:
\[
r^2 + \lambda^2 = 0
\implies r = \pm i \lambda
\]
Thus, the general solution for $X(x)$ is:
\[
X(x) = C_1 \cos(\lambda x) + C_2 \sin(\lambda x)
\]
Checking the boundary condition $u(0,t) = 0$:
\[
u(0,t) = X(0)T(t) = C_1 T(t) = 0
\implies C_1 = 0 \implies X(x) = C_2 \sin(\lambda x)
\]
Now checking the second boundary condition $u(L,t) = 0$:
\[
u(L,t) = X(L)T(t) = C_2 \sin(\lambda L) T(t) = 0
\implies \sin(\lambda L) = 0
\]
From the unit circle, $\sin (z) = 0$ when $z = n\pi$ for some integer $n$. Thus, we have:
\[
\lambda L = n \pi \implies \lambda = \frac{n \pi}{L}, \quad n \in \NN
\]
Thus, the values of $k$ that yield non-trivial solutions are:
\[
k_n = -\lambda_n^2 = -\left(\frac{n \pi}{L}\right)^2, \quad n \in \NN
\]
And the solution for $X(x)$ is:
\[
X_n(x) = \sin\left(\frac{n \pi x}{L}\right), \quad n \in \NN
\]
These individual $X_n$ are called the eigenfunctions of the problem, and the corresponding $k_n$ are called the eigenvalues.

\subsubsection{Finding $T(t)$}\label{subsubsec.finding_T}
\[
T' - a^2 k T = 0
\implies T' + a^2 \lambda^2 T = 0
\]
This is a first-order linear ODE we can solve via separation:
\[
\diff{T}{t} = -a^2 \lambda^2 T
\implies \frac{1}{T} \diff{T}{t} = -a^2 \lambda^2
\implies \int \frac{1}{T} dT = -a^2 \lambda^2 \int dt
\]
\[
\implies \ln |T| = -a^2 \lambda^2 t + C
\implies T(t) = A e^{-a^2 \lambda^2 t}
\]

\subsection{Constructing the General Solution}\label{subsec.general_solution}
From the previous sections, we have found that the eigenfunctions of the problem are $X_n(x) = \sin\left(\frac{n \pi x}{L}\right)$ and the corresponding time functions are $T_n(t) = e^{-a^2 \lambda_n^2 t}$. Thus, the product of these gives us a family of solutions to the heat equation:
\[
u_n(x,t) = X_n(x) T_n(t) = \sin\left(\frac{n \pi x}{L}\right) e^{-a^2 \left(\frac{n \pi}{L}\right)^2 t}, \quad n \in \NN
\]
Combining these solutions, we can construct the general solution to the heat equation as an infinite sum of these eigenfunctions:
\begin{equation} \label{eqn.heat_eqn_sol}
    u(x,t) = \sum_{n=1}^\infty b_n \sin\left(\frac{n \pi x}{L}\right) e^{-a^2 \left(\frac{n \pi}{L}\right)^2 t}
\end{equation}
Checking the final initial condition:
\[
u_n(x,0) = \sin\left(\frac{n \pi x}{L}\right) = f(x)
\]
depending on the initial condition $f(x)$, this could be just a matter of choosing the right $n$ to match $f(x)$. But what if $f$ is not a trigonometric function? Can we represent an arbitrary $f$ as a sum of trigonometric functions? Enter, \textbf{Fourier}! This is how we can find the coefficients $b_n$ in the general solution \eqref{eqn.heat_eqn_sol}. We do not have the tools yet to do this, but we can after a brief introduction to Fourier series in the next section.

After the next section we will be able to show that the coefficients $b_n$ can be found by the following formula:
\[
b_n = \frac{2}{L} \int_0^L f(x) \sin\left(\frac{n \pi x}{L}\right) dx, \quad n \in \NN
\]

This now gives us a complete solution to the heat equation \eqref{eqn.heat_eqn} with the given initial and boundary conditions.
\begin{equation}\label{eqn.heat_eqn_complete_sol}
    u(x,t) = \sum_{n=1}^\infty \left( \frac{2}{L} \int_0^L f(x) \sin\left(\frac{n \pi x}{L}\right) dx \right) \sin\left(\frac{n \pi x}{L}\right) e^{-a^2 \left(\frac{n \pi}{L}\right)^2 t}
\end{equation}